\section{Limitations}
\label{sec:limitations}

\paragraph{The target and its labels.} This is the deepest limitation. In our main design each
fold's target pool is built from the graded lenses of the other folds and an equal number of
random galaxies, and every target batch is balanced 1:1 using the catalogue labels
(Section~\ref{sec:labels}). A search on new survey data has no graded lenses to build such a
pool from; it would face the survey's own rarity, about one lens found per 2200 objects
searched~\citep{lines2025simstosky}. Our dilution experiment
removes the balancing of the batches but keeps a pool built from catalogue lenses, and its
1:\dlRatio{} is still about \dlDenserWord{} times denser in lenses than the survey's 1:2200; we
tested no target at survey prevalence, so our results do not transfer unaltered to that
setting. The dilution changes more than the proportion (Section~\ref{sec:dilute-methods}). It ran one ratio, with one seed per fold, and CORAL's diluted run is entangled
with the head its frozen stage produced (Section~\ref{sec:results-dilute}). The theory that
motivated balancing covers marginal alignment, which is CORAL and
DANN here~\citep{zhao2019invariant, zhao2020gls}, and a benchmark finds such methods fragile
under shifts in class proportions~\citep{classimbalanceda}; it does not speak directly to FreeMatch or
BNM, and we applied the balanced target to every arm to keep the comparison matched. The
closest precedent shares the constraint: its balanced target of 1754 lenses and 1754
non-lenses is also built from known lens candidates~\citep{parul2024}.

Balancing is not the only remedy for differing class proportions. \citet{zhao2020gls} estimate
the class ratio between domains from the source confusion matrix and the target predictions,
then reweight the adversarial and classification losses by it, which needs no target label. We
tried one such method, IW-CDAN, only in early frozen runs with an imbalanced target, where it
collapsed together with DANN (Section~\ref{sec:labels}). Label-shift estimation methods that
need no target labels, such as re-estimating the class prior by
expectation--maximisation~\citep{saerens2002} or black-box shift
estimation~\citep{lipton2018}, are a natural next step.

\paragraph{What the target signal contains.} Because every target batch of the main design is
half lenses and half galaxies, the adaptation arms see not only real images but a known class
proportion, and each
objective can use it: BNM's diversity term favours a batch whose predictions split evenly,
FreeMatch's fairness term targets a uniform class marginal, and DANN and CORAL align the
target's feature distribution with that of the balanced source. Learning from groups of
unlabelled examples whose class proportions are known has been studied as a learning problem in
its own right~\citep{quadrianto2009}. The $\lambda = 0$ control removes both the target images'
gradient and this proportion, apart from the few half-precision steps it skips
(Section~\ref{sec:arms}), so our contrasts cannot say how much of the gain comes from
adapting to real images and how much from the known proportion. The dilution experiment answers
this only in part. DANN keeps a detectable gain over the control without batches balanced 1:1
(Section~\ref{sec:results-dilute}), although the comparison with the balanced target cannot
rule out that it lost much of its gain. Every objective still sees a target whose composition
we fixed with the labels, and FreeMatch's fairness term still assumes a uniform class mix, so
we still cannot separate adapting to real images from the class proportion.

\paragraph{The diluted runs' numerics.} CORAL's unfrozen stage on the diluted target skipped its
first \dlCoralSkipMin{}--\dlCoralSkipMax{} optimiser updates on every fold, before its loss
scaling settled, and none after, so it applied at least \dlCoralAppliedPct{}\% of its budget's
updates. Its inherited
feature-norm guard, which our pre-registration made advisory, fired on every step. Its first
attempt on fold 1 stopped at step 1 on that guard, an implementation error that produced no
data, and the third part of our pre-registration permitted the rerun
(Appendices~\ref{sec:prereg-dilute} and~\ref{app:dilute10}).

\paragraph{A weak simulation-only baseline.} Our simulation-only baseline recovers
$\dLrBaseline{}$\% of the benchmark's completeness, so the
gains we report are measured from a low starting point, against a control that is itself one
source-only fine-tune (Section~\ref{sec:results-ablation}).

\paragraph{When the control was added.} The control and its contrasts were added after the
adaptation arms' results were known; its own runs were pre-registered, but the decision to make
those comparisons was not (Section~\ref{sec:prereg}).

\paragraph{Fold dependence, four folds and one seed.} Every interval rests on four folds whose
target pools overlap (Section~\ref{sec:intervals}). Every reported result comes from a single
run per fold with one seed, so seed variance is not separated from fold variance anywhere. The
ranking analysis of Section~\ref{sec:resolution} is approximate for the same reason.

\paragraph{Test negatives and depth.} \FracFaintTest{}\% of our test negatives are fainter
than the benchmark's limit of
$I_E = 22.5$ (Section~\ref{sec:splits}). Some may be lenses: the exact-ID audit found
catalogued lenses in the original pool, and $\rvTopNegC{}$ of the highest-scoring negatives in
the cleaned pool are grade-C
candidates (Section~\ref{sec:findmiss}). Restricting the pool to the benchmark's depth leaves
every contrast we recomputed (Table~\ref{tab:bright}) with the same zero-inclusion
(Appendix~\ref{app:bright}), but purity values are
not directly comparable with the benchmark's. Our positives were found mostly by
machine-learning finders, one of them built on Zoobot, the encoder we use, so recall on them may
overstate recall on lenses such finders miss~\citep{lines2025discoveryengineC, discoveryengineA}.

\paragraph{The simulations.} Our lensed sources are parametric profiles fitted to COSMOS
galaxies, so they lack the clumps and irregular structure of real high-redshift sources; in HST
imaging only about a fifth of galaxies are fully described by such
models~\citep{mandelbaum2014great3}. The deflector velocity dispersions follow a local
($z \lesssim 0.1$) velocity dispersion function, while our deflectors lie at a median redshift
of \simZlMedian{}, and we apply no evolution. \DeflUnphysN{} deflectors carry catalogue values outside the
physical range ($\log M_* > \DeflUnphysMass{}$ or an absolute VIS magnitude below $\DeflUnphysMag{}$),
and \DeflUnphysDrawn{} of these host \DeflUnphysLenses{} of the simulated lenses.

\paragraph{Software.} We did not pin or record the versions of \texttt{timm} and GalSim used in
training and simulation (Appendix~\ref{app:repro}).
